% ============================================================================
%  CONCLUSIONES Y RECOMENDACIONES
% ============================================================================

\section{Conclusiones del Ejercicio de Auditoría}

El desarrollo integral de la auditoría de seguridad y escaneo de vulnerabilidades sobre la infraestructura del servidor objetivo (\texttt{192.168.100.85}) ha permitido consolidar conclusiones técnicas y operativas de alto valor académico y profesional:

\begin{enumerate}
  \item \textbf{Precisión del Análisis Ofensivo con Evasión Táctica:} La auditoría evidenció que los escáneres automatizados convencionales pueden ser neutralizados o degradados por mecanismos de retención (\textit{tarpitting}) y límites de concurrencia (\textit{rate-limiting}) implementados en Nginx y el firewall. No obstante, la aplicación de técnicas tácticas avanzadas como el escaneo con señuelos aleatorios (\textit{Decoy Scanning} \texttt{-D RND:20}), la exclusión de sondas DoS y la restricción de tiempo por script (\texttt{--script-timeout 30s}) permitieron completar el escaneo de los 16 servicios en tan solo 49.66 segundos, demostrando que la seguridad por ofuscación o retardo es insuficiente ante auditores experimentados.
  
  \item \textbf{Prevalencia de Fallos de Configuración y Cifrado:} Los hallazgos de severidad alta identificados en la evaluación (el intercambio anónimo Diffie-Hellman en Postfix SMTP, la exposición pública del puerto PostgreSQL en \texttt{0.0.0.0:5432} y la persistencia de verbos de depuración \texttt{DEBUG} en el servicio de Webmail) no se originaron en fallos de día cero (\textit{0-days}) del código base, sino en omisiones de endurecimiento (\textit{hardening}) y uso de parámetros permisivos por omisión. Asimismo, mediante la inspección forense de tráfico con Wireshark, se corroboró empíricamente la circulación de credenciales y cookies de sesión en texto plano en aplicaciones administrativas desprovistas de encapsulamiento TLS.

  \item \textbf{Efectividad Comprobada del Monitoreo y Prevención (Suricata e IDS/IPS):} La integración del motor Suricata junto con la interfaz gráfica EveBox demostró una capacidad sobresaliente para decodificar flujos de red en tiempo real e interceptar sondas ofensivas de inyección SQL y explotación web. Esto ratifica que la combinación de análisis de firmas y visibilidad gráfica centralizada reduce drásticamente el tiempo medio de detección e investigación de incidentes.

  \item \textbf{Complementariedad de Herramientas DAST y de Red:} El uso conjunto de Nmap para la auditoría de puertos de transporte y OWASP ZAP para el escaneo de aplicaciones web (como se demostró sobre el portal Uyuni Way en el puerto 8085) permitió una cobertura integral tanto a nivel de infraestructura como a nivel de lógica de aplicación web, visibilizando alertas que escáneres de puertos puros no logran detectar.

  \item \textbf{Efectividad y Viabilidad de las Medidas de Mitigación y Endurecimiento:} La aplicación del plan de mitigación técnica propuesto —restringiendo el acceso a PostgreSQL a la interfaz local (\texttt{127.0.0.1}), limitando el procesamiento de rangos en Nginx (\texttt{max\_ranges 1;}), excluyendo suites criptográficas anónimas (\textit{aNULL}) en Postfix e imponiendo parámetros Diffie-Hellman robustos de 2048 bits— neutraliza de manera definitiva los vectores de explotación más críticos sin comprometer la continuidad de las operaciones del servidor.

  \item \textbf{Consolidación de Identidad Digital y No Repudio (PAdES ISO 32000-2):} La incorporación y evaluación de la Plataforma de Firma Digital Institucional demostró la factibilidad técnica de integrar servicios criptográficos avanzados (emisión PKI automatizada, certificados X.509 v3, algoritmos SHA-256/RSA-4096 y firmas PAdES incrementales) dentro de una infraestructura Linux contenerizada. La aplicación rigurosa de políticas Zero Trust (Auth Gate) y el confinamiento del backend tras un proxy inverso TLS aseguran que las garantías probatorias de autenticidad, integridad y no repudio se preserven sin incrementar la superficie vulnerable del servidor.
\end{enumerate}

\section{Recomendaciones para la Infraestructura y Práctica Operativa}

A partir de los resultados obtenidos y de las dificultades operativas experimentadas durante el despliegue del laboratorio, se formulan las siguientes recomendaciones estratégicas y técnicas:

\begin{enumerate}
  \item \textbf{Superación de Limitaciones de Escenario y Adquisición de Equipamiento de Red:} Preparar un entorno de pruebas fidedigno para auditar un proyecto propio presenta una alta complejidad técnica cuando se carece de equipamiento físico dedicado de conectividad y seguridad (tales como routers corporativos, switches administrables con soporte para VLANs y firewalls de hardware dedicados). La ausencia de estos dispositivos impide la implementación de una topología de red verdaderamente funcional y segmentada (con separación estricta entre zonas DMZ, Intranet, Gestión y Base de Datos). Esto deriva en una pérdida sensible de la práctica en el diseño de arquitecturas seguras y en la formulación de políticas de filtrado avanzadas en hardware perimetral, obligando a simular controles en entornos planos donde el tráfico de prueba interfiere con los servicios evaluados.
  
  \item \textbf{Segmentación y Aislamiento de Red (Implementación de DMZ):} Para solventar la exposición directa de servicios críticos, se recomienda configurar interfaces virtuales (VLANs 802.1Q) o subredes Docker internas sin exposición al exterior (\textit{internal networks}). Servicios de base de datos como PostgreSQL (5432) y herramientas administrativas como Portainer (8000) deben escuchar exclusivamente en interfaces de bucle local (\texttt{127.0.0.1}) o requerir túneles SSH/VPN para su administración.

  \item \textbf{Endurecimiento Criptográfico de Protocolos de Correo y Web:} Desactivar de forma perentoria las suites de cifrado débiles y anónimas (\texttt{!aNULL:!eNULL:!ADH}) en la configuración de Postfix (\texttt{main.cf}), aplicar grupos Diffie-Hellman robustos de 2048 bits generados localmente y mitigar el vector de agotamiento de memoria Apache/Nginx ByteRange mediante la directiva \texttt{max\_ranges 1;}.

  \item \textbf{Auditorías Continuas con DAST y Automatización DevSecOps:} Incorporar análisis dinámicos programados con OWASP ZAP en modo demonio dentro del ciclo de desarrollo web, resolviendo las alertas de severidad media (ausencia de cabeceras de protección CSP, HSTS y fuga de métodos HTTP no requeridos) previo al paso a producción.

  \item \textbf{Auditoría y Gestión de Alertas con EveBox y Suricata:} Establecer un procedimiento sistemático de revisión periódica del archivo \texttt{eve.json} mediante EveBox, creando alertas automáticas hacia el canal de administración del equipo ante la detección de anomalías de red o intentos de evasión de firewall.

  \item \textbf{Gestión del Ciclo de Vida de Claves y Rotación PKI:} Para la infraestructura de firma digital, se recomienda establecer un procedimiento formal de rotación periódica de la CA Raíz institucional y de las llaves maestras Fernet en reposo, complementado con la publicación de listas de revocación de certificados (CRL) y protocolos OCSP para la invalidación inmediata de credenciales comprometidas.
\end{enumerate}
